In this paper, we first provide a comprehensive review of the current advancements in the field of Graph Unlearning, highlighting its practical applications and systematically categorizing existing algorithms based on their technical characteristics. Building on this foundation, we introduce OpenGU, the first unified and comprehensive benchmark for GU, which integrates 16 state-of-the-art GU algorithms and 37 datasets across multiple domains. OpenGU further extends the flexibility of GU algorithms, allowing seamless combinations across three downstream tasks and three types of unlearning requests. Using the standardized and fair evaluation framework provided by OpenGU, we conduct extensive experiments to assess the advancements of GU methods from the perspectives of effectiveness, efficiency, and robustness. These experiments not only reveal the notable breakthroughs in the field but also expose critical limitations in existing approaches. To inspire further research, we outline the major challenges currently faced by GU and propose promising directions for future exploration.

\textbf{Designing Generalized GU Frameworks for Diverse Tasks} (\textit{C1} and \textit{C2}). While some existing methods demonstrate strong predictive performance for specific tasks, their underlying principles often make it difficult to adapt to varying downstream tasks and unlearning requests. Furthermore, the current design of unlearning requests remains relatively simplistic, whereas real-world applications often require handling mixed unlearning requests or subgraph-level unlearning. Achieving consistently superior predictive capabilities across such complex and varied scenarios remains a considerable challenge. Therefore, developing more generalized frameworks becomes imperative.

\textbf{Unified Metrics for Evaluating Forgetting} (\textit{C3} and \textit{C4}). The current methods for assessing the forgetting capability of GU approaches remain insufficient and leave significant gaps to be addressed. These methods are often tightly coupled with specific unlearning requests and downstream tasks, making them less effective in handling diverse combinations of scenarios. Moreover, determining whether the information targeted for removal has been thoroughly unlearned from theoretical perspective remains a critical yet underexplored challenge. Furthermore, GU algorithms must strike a balance between reasoning and forgetting. Future research should move beyond the current paradigm of independently assessing these two aspects, striving instead for a unified metric that evaluates models from an integrated perspective, ensuring a comprehensive understanding of their capabilities.

\textbf{Enhancing Algorithm Efficiency} (\textit{C5} and \textit{C6}). While theoretical analysis provides valuable insights, the practical performance of current GU methods often falls short, especially when scaling to large datasets with millions of nodes. These methods commonly encounter OOT or OOM issues. To enable GU's effective deployment in large-scale scenarios, algorithms must be optimized for both efficiency and scalability, ensuring they can handle the demands of real-world data while avoiding these performance bottlenecks.

\textbf{Addressing Realistic Scenarios} (\textit{C7} and \textit{C8}). In practical applications, the presence of noise and incomplete datasets is an unavoidable challenge. However, current GU algorithms lack sufficient exploration and adaptation to such scenarios. Experimental results highlight significant weaknesses when dealing with noise and sparsity, particularly in terms of label and feature robustness. Future research should aim to broaden the scope of investigation, extending robustness analysis to encompass a wider range of real-world challenges and data imperfections.

OpenGU embodies the collective efforts and expertise of the current GU field, offering a comprehensive platform for research and development. As we move forward, we will continue to expand and enhance OpenGU, strengthening its support for future advancements in this area. Finally, we welcome feedback and encourage suggestions to refine our benchmark, further improving its effectiveness and user experience.